\section{A nearest-square guard: 28 unrestricted integers at degree eight}
The two independent positivity guards can be compressed without changing any
Pell equation, compiler coefficient, or arithmetic register.
For integer $x,y$, put
\[
 G(x,y,t)=xy-(t^2-x-y+1)^2.
\]
Then $x,y>0$ if and only if $G(x,y,t)>0$ for some integer $t$.
Indeed, $G>0$ forces $xy>0$, while $x=-a,y=-b$ with $a,b>0$
would give $(t^2+a+b+1)^2>ab$. Conversely, let $N=x+y-1$.
Then $xy\geq N\geq1$. Writing $m=\lfloor\sqrt N\rfloor$, choose
$t=m$ if $N-m^2\leq m$ and $t=m+1$ otherwise. The square error
has absolute value at most $m$, and its square is strictly less than $N$
(the exact-square case has error zero). Thus $G>0$.

Every positive integer is
$\operatorname{Pos}(p,q,r)=1+p^2+q^2+r(r+1)$ by the three-square theorem
applied to $4n-3$. Therefore the single residual
\[
 xy-(t^2-x-y+1)^2-1-p^2-q^2-r(r+1)
\]
replaces two separate positive ternary forms, using four private integer
coordinates rather than six. In the reviewed 30-integer octic, take
$x=F\alpha$ and $y=F(\beta J+1-F-\alpha)$.
Both are quadratic, so the new residual is quartic. All other residuals
remain unchanged, including the guard on $\eta\zeta$.
The resulting sum of 14 squared residuals has 28 unrestricted integer
variables and degree exactly eight; the unchanged coefficient
$[a^4f^4]=1$ certifies exactness.

The equivalence is existential, not a bijection of witness tuples. From
an old zero choose the nearest-square coordinate and a positive ternary
representation of $G$. From a new zero restore independent positive
ternary representations of $x$ and $y$. Composing these maps with the
pinned source theorem gives the ZFC-provable equivalence between absence
of integer zeros and $\operatorname{Con}(\mathrm{ZFC})$.
No global minimality, external priority, or new proof-assistant certificate
is asserted.

A second application gives a higher-degree companion. With
$z=\eta(1-V+jc+hE-\eta)$ and $g=G(x,y,t)$, replace all three old
positive ternary guards by
\[
 gz-(u^2-g-z+1)^2=\operatorname{Pos}(p,q,r).
\]
Applying the proved lemma twice shows that this is equivalent to
$x,y,z>0$. Nine private coordinates become five, leaving 26 integers
and 13 residuals. The new residual has coefficient $-1$ at $t^8$,
so its square has coefficient $1$ at $t^{16}$; the other residuals do
not involve $t$. The exact total degree is therefore sixteen.
